We demonstrate that the odd values $\zeta(2k+1)$ of the Riemann zeta
function are structurally determined by $\varphi=(1+\sqrt{5})/2$ and~$\pi$.
The Euler product of $\zeta$ is built from local factors $f_p(s)$, one
for each prime.  We prove that the Frobenius lift at each prime $p\ne5$,
congruent modulo $p$ to $\golden^p = F_p\golden + F_{p-1}$---where $F_n$ is
the $n$-th Fibonacci number---determines the splitting type of $p$ in
$\mathbb{Z}[\golden]$, and hence, with the ramified prime $5$, which admits
no lift, every $f_p(s)$. These lifts constitute $\mathbb{F}_1$-descent data
in the sense of Borger away from $5$. Through the Dedekind factorisation
\[
  \zeta(s) = \frac{\zeta_{\mathbb{Q}(\sqrt{5})}(s)}{L(s,\chi_5)},
\]
this determines $\zeta(s)$ completely: the Frobenius structure implies
the Euler product at every level (primes, splitting types, local factors,
$L$-function), each implication proved, anchored by the base case
$L(1,\chi_5) = 2\log\golden/\sqrt{5}$
(the class-number formula for $\mathbb{Q}(\sqrt{5})$),
which expresses the $L$-function value entirely in terms of~$\golden$;
and the factorisation is a natural isomorphism over $\mathbb{Q}$ and the
quadratic fields of discriminant dividing~$20$.
The determination fixes the odd values through the pentagonal arithmetic
without evaluating them, and by itself it does not decide their
arithmetic nature.  The arithmetic is formalised in Lean 4 with Mathlib, and each
statement carries its tier: proved, assumed from the literature, numerical, or open.
